\section{Описание практической работы}
\label{sec:work} \index{work}

\subsection {Исправление генерации отладочной информации в компиляторе LLVM}
\label{work:clang_fix}

В процессе проведения различных экспериментов с генерацией отладочной информации,
было обнаружено, что для кортежных типов, таких как \texttt{vint8m4x2\_t}, компилятор LLVM
генерирует некорректную отладочную информацию о размере таких типов.
Компилятор не учитывал множитель \texttt{NFIELDS}.

Ниже приведена часть исходного кода LLVM с исправлением этой ошибки:

\begin{samepage}
\begin{lstlisting}[language=C++, caption=Исправление генерации информации о размере некоторых типов, label=lst:code:llvm_fix]
bool Fractional = false;
unsigned LMUL;
unsigned NFIELDS = Info.NumVectors;
unsigned FixedSize = ElementCount * SEW;
if (Info.ElementType == CGM.getContext().BoolTy) {
    // Mask type only occupies one vector register.
    LMUL = 1;
} else if (FixedSize < 64) {
    // In RVV scalable vector types, we encode 64 bits in the fixed part.
    Fractional = true;
    LMUL = 64 / FixedSize;
} else {
    LMUL = FixedSize / 64;
}

// Element count = (VLENB / SEW) x LMUL x NFIELDS
SmallVector<uint64_t, 12> Expr(
    // The DW_OP_bregx operation has two operands: a register which is
    // specified by an unsigned LEB128 number, followed by a signed LEB128
    // offset.
    {llvm::dwarf::DW_OP_bregx, // Read the contents of a register.
    4096 + 0xC22,             // RISC-V VLENB CSR register.
    0, // Offset for DW_OP_bregx. It is dummy here.
    llvm::dwarf::DW_OP_constu,
    SEW / 8, // SEW is in bits.
    llvm::dwarf::DW_OP_div, llvm::dwarf::DW_OP_constu, LMUL});
if (Fractional)
    Expr.push_back(llvm::dwarf::DW_OP_div);
else
    Expr.push_back(llvm::dwarf::DW_OP_mul);
// NFIELDS multiplier
if (NFIELDS > 1)
    Expr.append({llvm::dwarf::DW_OP_constu, NFIELDS, 
                 llvm::dwarf::DW_OP_mul});
// Element max index = count - 1
Expr.append({llvm::dwarf::DW_OP_constu, 1, 
             llvm::dwarf::DW_OP_minus});
\end{lstlisting}
\end{samepage}

Соответствующий патч с добавлением необходимого множителя \texttt{NFIELDS} был
мною предложен и принят в исходный код компилятора LLVM.

\subsection {Реализация в GDB}
\label{work:realisation}

Для того, чтобы поддержать RISC-V Vector Calling Convention в отладчике GDB, необходимо
было реализовать корректную передачу аргументов с векторными типами в функцию перед ее вызовом. 
Вторым требованием
является сохранение значений определенных регистров и их восстановление после окончание
исполнения функции, однако GDB самостоятельно сохраняет значения всех регистров перед вызовом
и восстанавливает их после, поэтому реализовывать данное требование не потребовалось.

Процесс размещения аргументов функции можно условно разделить на две стадии: определение корректного
местоположения аргументов функции в соответствии с ABI и размещение этих аргументов. В рамках
решения поставленных задач был модифицирован процесс установление правильного размещения
аргументов функции, добавлена поддержка определения правильного местоположения для векторных
аргументов.

Итак, в первую очередь, для определения векторного типа по его имени была написана следующая простая функция:

\begin{samepage}
\begin{lstlisting}[language=C++, caption=Функция для определения векторных типов, label=lst:func:is_rvv_type]
static bool
is_rvv_type (struct type *type)
{
  if (type)
  {
    if (type->name () && 
        (strncmp (type->name (), "__rvv", 5) == 0))
      return true;
    else if (type->code () == TYPE_CODE_TYPEDEF)
      return is_rvv_type (type->target_type ());
  }

  return false;
}
\end{lstlisting}
\end{samepage}

Например, тип \texttt{\_\_rvv\_int8m1\_t} (\texttt{vint8m1\_t}).

Для выбора способа передачи переменной векторного типа (через регистры или через память) была написана функция:

\begin{samepage}
\begin{lstlisting}[language=C++, caption=Функция для определения способа передачи значения аргумента, 
	           label=lst:func:riscv_call_arg_vector, breaklines=false]
static void
riscv_call_arg_vector (struct riscv_arg_info *ainfo,
                       struct riscv_call_info *cinfo,
                       struct type *func_arg_type)
{
  gdb_assert (((ainfo->length % cinfo->vlenb) == 0)
                || ainfo->length < cinfo->vlenb);

  if (!riscv_assign_vec_reg_location (ainfo, &cinfo->vector_regs,
                                      func_arg_type, cinfo->vlenb))
    {
      // Try to pass value by reference
      ainfo->argloc[0].loc_type = riscv_arg_info::location::by_ref;
      cinfo->memory.ref_offset = align_up (cinfo->memory.ref_offset, 
                                           ainfo->align);
      ainfo->argloc[0].loc_data.offset = cinfo->memory.ref_offset;
      cinfo->memory.ref_offset += ainfo->length;
      ainfo->argloc[0].c_length = ainfo->length;

      if (!riscv_assign_reg_location (&ainfo->argloc[1], 
                                      &cinfo->int_regs, 
                                      cinfo->xlen, 0))
        riscv_assign_stack_location (&ainfo->argloc[1], 
                                     &cinfo->memory,
                                     cinfo->xlen, cinfo->xlen);
    }
}
\end{lstlisting}
\end{samepage}

В первую очередь пытаемся разместить значение на векторных регистрах (в соответствии с RISC-V Vector
Calling Convention) с помощью функции \texttt{riscv\_assign\_vec\_reg\_location} (подробное описание
функции будет ниже). В случае неудачи, пытаемся передать значение по адресу, который, в свою очередь,
пытаемся сначала разместить на регистре общего назначения, используя  \texttt{riscv\_assign\_reg\_location}, 
или, на стеке, с помощью \texttt{riscv\_assign\_stack\_location}.

Реализация функции \texttt{riscv\_assign\_vec\_reg\_location}:

\begin{samepage}
\begin{lstlisting}[language=C++, caption=Функция для выбора набора регистров для передачи аргумента, 
			   label=lst:func:riscv_assign_vec_reg_location, breaklines=false]
static bool
riscv_assign_vec_reg_location (struct riscv_arg_info *ainfo,
                               struct riscv_vector_arg_reg *reg,
                               struct type *func_arg_type, int vlenb)
{
  /* Amount of data that is stored on one register */
  int data_length = vlenb;   
  
  struct riscv_arg_info::location *loc0 = &ainfo->argloc[0];
  struct riscv_arg_info::location *loc1 = &ainfo->argloc[1];
   
  const char *func_arg_type_name
    = (func_arg_type) ? func_arg_type->name () : nullptr;
   
  /* Try to get info about LMUL*/
  const char *lmul_pos = (func_arg_type_name) ? 
                          strstr (func_arg_type_name, "m") : nullptr;
  bool is_fractional_lmul = (lmul_pos && 
                             strncmp (lmul_pos, "mf", 2) == 0) ? 
                             true : false;
  int lmul = (lmul_pos) ? 
              ((is_fractional_lmul) ? 
                lmul_pos[2] - '0' : lmul_pos[1] - '0') : 1;
   
  /* If it is a segment type, it's name should ends with `xN_t`, 
     where N is * NFIELDS*/
  const char *nfields_pos = (func_arg_type_name) ? 
                          strstr (func_arg_type_name, "x") : nullptr;
  int nfields = (nfields_pos) ? nfields_pos[1] - '0' : 1;
   
  int num_required_regs = (ainfo->length <= vlenb) ? 
                           1 : ainfo->length / vlenb;
   
  if (is_fractional_lmul)
    {
      num_required_regs = nfields;
      data_length = data_length / lmul;
    }
   
  int first_regnum = -1;
   
  /* Representation of vector segmented types, like vint32m4x2_t,
      is different in Clang and GCC.
      In Clang, all types are arrays.
      In GCC, segmented types are struct's, 
      other vector types are arrays. */
  if (ainfo->type->code () == TYPE_CODE_ARRAY)
    {
      if (ainfo->type->target_type ()->code () == TYPE_CODE_BOOL)
        first_regnum = reg->try_use_v0 ();
      else
        first_regnum = reg->get_interval_start (num_required_regs, 
                                                nfields);
    }
  else if (ainfo->type->code () == TYPE_CODE_STRUCT)
    {
      first_regnum = reg->get_interval_start (num_required_regs, 
                                              nfields);
    }
   
  int last_regnum = first_regnum + num_required_regs - 1;
   
  if (first_regnum != -1)
    {
      loc0->loc_type = riscv_arg_info::location::in_several_regs;
      loc0->loc_data.regno = first_regnum;
      loc0->c_length = data_length;
      loc0->c_offset = 0;

      loc1->loc_type = riscv_arg_info::location::in_several_regs;
      loc1->loc_data.regno = last_regnum;
      loc1->c_length = 0;
      loc1->c_offset = 0;
   
      return true;
    }
   
  return false;
}
\end{lstlisting}
\end{samepage}

В данной функции мы сначала определяем параметры LMUL и NFIELDS. К сожалению, эти параметры
пока невозможно получить более \enquote{чистым} способом -- из отладочной информации, т.~к. их там
попросту нет. Мы обладаем только информацией о размере типа, из которой мы можем сделать
вывод о количестве необходимых векторных регистров, но не о множителе LMUL, необходимом
для корректной реализации векторного соглашения о вызовах.
Далее, вызывая
метод \texttt{get\_interval\_start} у объекта \texttt{struct riscv\_vector\_arg\_reg *reg},
который сохраняет в себе информацию об уже занятых регистрах, мы получаем номер первого
регистра в наборе (и вычисляем номер последнего регистра), в котором будет сохранено 
необходимое значение.

Как было сказано выше, для сохранения информации об уже занятых регистрах, была
написана структура данных \texttt{riscv\_vector\_arg\_reg}:

\begin{samepage}
\begin{lstlisting}[language=C++, caption=Структура для хранения информации о свободных векторных регистрах, 
			label=lst:struct:riscv_vector_arg_reg, breaklines=false]
struct riscv_vector_arg_reg_interval
{
  int first;
  int last;
};

struct riscv_vector_arg_reg
{
  riscv_vector_arg_reg (int first, int last)
      : available_intervals{ { first, last } }
  {
    /* Nothing. */
  }

  int
  get_interval_start (int count, int NFIELDS = 1)
  {
    gdb_assert (count % NFIELDS == 0);

    int LMUL = count / NFIELDS;
    for (auto cur_interval = available_intervals.begin ();
         cur_interval != available_intervals.end (); ++cur_interval)
      {
        int cur_first = cur_interval->first;
        int cur_last = cur_interval->last;
        if ((cur_first - RISCV_V0_REGNUM) % LMUL != 0) 
             /* According to RISC-V Vector ABI,
             first register number should be
             a multiple of LMUL */
          {
            cur_first -= (cur_first - RISCV_V0_REGNUM) % LMUL;
            cur_first += LMUL;
          }

        if ((cur_last - cur_first + 1) >= count)
          {
          if (cur_first > cur_interval->first)
            {
              available_intervals.insert 
                (cur_interval, { cur_interval->first, cur_first - 1 });
            }

            cur_interval->first = cur_first + count;

            return cur_first;
	  }
      }

    return -1;
  }

  int
  try_use_v0 ()
  {
    if (is_v0_used)
      return get_interval_start (1);
    else
      {
        is_v0_used = true;
        return RISCV_V0_REGNUM;
      }
  }

  bool is_v0_used
      = false; /* To pass masks we should use v0, if it's available */

  std::list<riscv_vector_arg_reg_interval> available_intervals;
};
\end{lstlisting}
\end{samepage}

Данные о незанятых векторных регистров хранятся в виде интервалов вида 
\texttt{\{first, last\}} с номерами регистров. В случае, если нам необходимо
выделить какой-то набор регистров, лежащий в интервале, этот интервал
разбивается на 3 части: необходимый нам набор, интервал до и интервал после
необходимого набора. Эти новые интервалы сохраняются взамен старого интервала\
(в случае, если они не являются пустыми).

Кроме того, отдельно отслеживается занятость регистра \texttt{v0}, т.~к. он
используется для передачи векторных масок.

Последняя важная часть, это обработка введенного варианта расположения
переменной на регистре \texttt{riscv\_arg\_info::location::in\_several\_regs}
в функции \texttt{riscv\_push\_dummy\_call}:

\begin{samepage}
\begin{lstlisting}[language=C++, caption=Запись необходимого значения в соответствующие регистры в кэше, 
		label=lst:switch:in_several_regs, breaklines=false]
case riscv_arg_info::location::in_several_regs:
{
  const gdb_byte *cur_contents = info->contents;
  int cur_c_length = info->argloc[0].c_length;
  for (int cur_regno = info->argloc[0].loc_data.regno;
       cur_regno <= info->argloc[1].loc_data.regno; cur_regno++)
    {
      riscv_regcache_cooked_write (cur_regno, cur_contents,
                                   cur_c_length, regcache,
                                   call_info.flen);
      cur_contents += cur_c_length;
    }

  second_arg_length = 0;
}
\end{lstlisting}
\end{samepage}







\newpage